\begingroup
\sloppy
\section{Unified Mobile Memory Governance}
\label{sec:pilot-unified-mobile-memory}

The unified Pilot phone application now exposes the existing AION private-memory authority as a
user-controlled Personal surface.  Memory is not treated as an invisible model prompt or copied
into a separate mobile database.  The mother brain remains the source of truth, while the trusted
phone provides inspection, explicit entry, correction, scope control, export and deletion.

\subsection{Identity and Guardian Boundary}

Every request is signed by the phone's non-exportable Ed25519 key and validated against its active
mother-issued certificate and capability lease.  Three independent scopes are used:
\texttt{memory.read}, \texttt{memory.write} and \texttt{memory.export}.  The private persona in
the certificate is always the requester.  A requester may target only themselves or a child profile
whose active guardian is that requester.  Sibling, unrelated child and other-adult memory access
fails closed before the memory store is read.

The interface identifies whose memory is open and marks a child surface as guardian managed.  It
does not merge an adult's records with a child's records, and it does not make a child profile a
general route into household memory.

\subsection{Inspectable Records and Scope}

Each projected record contains a bounded kind, human-readable summary, explicit scope, creation or
correction time and integrity hash.  The supported scopes are \texttt{private} and
\texttt{household\_shared}.  New records default to private.  Moving a record into household scope
requires an affirmative confirmation on the trusted phone; it is never inferred from television
speech or silently widened by a model.

The person can correct the summary without replacing the record identity.  Pilot records the
correction time and recomputes its canonical integrity hash.  Create, correct, scope-change and
delete operations carry idempotency keys.  An identical retry returns the original result, whereas
reuse of the same key with different content is rejected.

\subsection{Export and Permanent Deletion}

Export produces a structured JSON package containing the selected profile, that profile's
memories, inspectable service consents and non-secret trusted-device metadata.  Recovery-code
hashes and device public keys are removed.  Export authority is separate from read authority so a
future organization or child policy can permit inspection without permitting bulk extraction.

Deletion requires a second, explicit irreversible confirmation in the phone interface and a signed
\texttt{memory.write} request.  The memory content is removed permanently.  The replay-safe
operation receipt remains so a retransmitted request cannot recreate ambiguity, but Pilot does not
claim the deleted content is recoverable.

\subsection{Mobile Experience and Verification}

The Personal shortcut rail now includes a dedicated Memory panel.  It provides an identity selector
for the current adult and their guardian-controlled children, an explicit-memory form, scope labels,
correction and scope controls, permanent deletion and a local JSON download.  Status messages state
whether a change was saved, shared, made private, exported or deleted.

Focused authority tests verify self access, guardian-child access, denial of other-adult access,
idempotent creation and deletion, correction and scope changes, and secret-free export.  Mobile
surface checks verify the signed endpoints and irreversible-delete wording.  The optimized Next.js
production build compiles successfully.  No live memory, account, provider or television state was
changed during automated verification.

\endgroup
